\chapter{System Evaluation}

\section{Introduction}

This chapter evaluates the system against the six objectives defined in
the Introduction. It presents quantitative results for both machine
learning components, discusses the behaviour of the end to end
application pipeline, and reflects critically on the limitations of the
approaches taken. Where results fall short of ideal, the reasons are examined honestly rather than being simplified or ignored, as understanding why a
component under performs is as important as knowing that it does.

The evaluation is structured in three parts. Sections
\ref{sec:eval-difficulty} and \ref{sec:eval-tags} cover the two
classifiers in detail, presenting training curves, model comparisons,
confusion matrices, and per class analysis. Section
\ref{sec:eval-application} evaluates the application level objectives:
the generation pipeline, the code execution engine, and the user
features. Section \ref{sec:eval-objectives} then maps the findings back
to each objective explicitly.

\section{Objective 3: Difficulty Classifier}
\label{sec:eval-difficulty}

\subsection{Dataset and Class Balance}

The difficulty classifier was trained on problems extracted from IBM
Project CodeNet \cite{puri2021codenet}, with difficulty labels derived
from acceptance rate percentile binning as described in the System
Design chapter. Because the labels were produced by dividing the sorted
acceptance rate distribution into three equal terciles, the resulting
class distribution is by construction approximately balanced, as shown
in Figure \ref{fig:difficulty-distribution}. Each class contains
approximately 1,280 problems, giving a total dataset of roughly 3,840
labelled problems. This balance is a direct consequence of the
percentile binning strategy and is an important positive property: it
means that a naive majority class baseline would achieve only around 33
percent accuracy, and it removes the need for class weighting or
oversampling techniques \cite{chawla2002smote} during training.

\begin{figure}[H]
    \centering
    \includegraphics[width=0.7\linewidth]{./images/evaluation/difficulty_distribution.png}
    \caption{Class distribution of the difficulty labels derived from
    acceptance rate percentile binning across the CodeNet dataset. The
    near equal counts across easy, medium, and hard confirm that the
    tercile binning strategy produced a balanced training set.}
    \label{fig:difficulty-distribution}
\end{figure}

\subsection{Model Comparison}

Seven candidate models were evaluated using five fold cross validation
\cite{kohavi1995crossval} scored on F1 macro: XGBoost
\cite{chen2016xgboost}, HistGradientBoosting, Gradient Boosting
\cite{friedman2001gradient}, Logistic Regression, Linear SVM
\cite{cortes1995svm}, Random Forest \cite{breiman2001randomforests},
and k Nearest Neighbours. Validation accuracy scores are shown in
Figure \ref{fig:difficulty-model-comparison}.

\begin{figure}[H]
    \centering
    \includegraphics[width=\linewidth]
        {./images/evaluation/model_comparison.png}
    \caption{Validation accuracy comparison across all seven candidate
    difficulty classifier models. XGBoost achieved the highest
    validation accuracy at 0.605, narrowly ahead of HistGradientBoosting
    at 0.597.}
    \label{fig:difficulty-model-comparison}
\end{figure}

XGBoost \cite{chen2016xgboost} achieved the highest validation accuracy
at 0.605, narrowly ahead of HistGradientBoosting at 0.597. Both
ensemble methods substantially outperformed the weaker baselines of k
Nearest Neighbours (0.469) and Linear SVM (0.544). The relatively small
spread between the top four models (0.543 to 0.605) suggests that the
difficulty classification task is intrinsically challenging, which is
expected given that the labels are derived from a proxy signal rather
than from ground truth human judgements \cite{zhou2017weaksupervision}.

\subsection{Final Model Results}

XGBoost was selected as the final model and evaluated on the held out
test set. The confusion matrix is shown in Figure
\ref{fig:difficulty-confusion-matrix}.

\begin{figure}[H]
    \centering
    \includegraphics[width=0.75\linewidth]{./images/evaluation/confusion_matrix.png}
    \caption{Confusion matrix for the final XGBoost difficulty
    classifier evaluated on the held out test set. Rows represent true
    labels and columns represent predicted labels.}
    \label{fig:difficulty-confusion-matrix}
\end{figure}

The per class results derived from the confusion matrix are summarised
in Table \ref{table:difficulty-results}.

\begin{table}[H]
    \centering
    \begin{tabular}{lcccc}
        \hline
        \textbf{Class} & \textbf{Correct} & \textbf{Total} &
        \textbf{Recall} & \textbf{Misclassified as} \\
        \hline
        Easy   & 175 & 260 & 0.67 & Medium (62), Hard (23) \\
        Medium & 125 & 259 & 0.48 & Easy (55), Hard (79) \\
        Hard   & 162 & 260 & 0.62 & Easy (21), Medium (77) \\
        \hline
        \textbf{Overall} & \textbf{462} & \textbf{779} &
        \multicolumn{2}{c}{\textbf{Accuracy: 0.593}} \\
        \hline
    \end{tabular}
    \caption{Per class performance of the XGBoost difficulty classifier
    on the test set, derived from the confusion matrix.}
    \label{table:difficulty-results}
\end{table}

The model performs best on the easy class (recall 0.67) and worst on
medium (recall 0.48). The medium class is consistently the hardest to
classify across all models, which is a well known phenomenon in ordinal
three class problems: the middle class shares boundary characteristics
with both neighbours and therefore receives the most misclassifications
in both directions. The hard class shows a notable asymmetry: 77 hard
problems are misclassified as medium but only 21 as easy, which makes
intuitive sense given that hard and medium problems are more similar
to each other than hard and easy problems are.

\subsection{Critical Analysis: The Proxy Label Problem}

The most significant limitation of the difficulty classifier is one
that cannot be resolved by changing the model: the labels themselves
are a proxy. Acceptance rate reflects how often competitive programmers,
ranging from beginners to experts, successfully solve a problem on a
judging platform. This is a reasonable but imperfect signal for
difficulty for several reasons. First, acceptance rate is influenced
by factors beyond intrinsic problem difficulty, including the
popularity of the problem (which determines how many attempts are
made), the typical skill level of the platform's user base at the time
the problem was active, and whether the problem has known edge cases
that trip up otherwise strong solutions. Second, the percentile
binning means that the label boundaries are arbitrary: a problem at
the 33rd percentile is labelled easy and one at the 34th is labelled
medium, even though their acceptance rates may differ by a fraction of
a percent.

A test set accuracy of approximately 0.59 on a balanced three class
problem represents a meaningful improvement over the random baseline
of 0.33 and the majority class baseline of 0.33. However, it also
means that roughly four in ten predictions are incorrect. In the
context of the platform, this is acceptable because the difficulty
label is presented to the user as a rough guide rather than a
guarantee. A question labelled easy that is genuinely medium does not
break the user experience. It merely provides slightly miscalibrated
guidance. The value of the component lies in having some automated
difficulty signal on AI generated content that would otherwise carry
no label at all.

\section{Objective 2: Topic Tag Classifier}
\label{sec:eval-tags}

\subsection{Dataset Analysis}

The topic tag classifier was trained on LeetCode \cite{leetcode2025}
problems as a multi label classification task, since a single problem
may be associated with multiple algorithmic topic tags. After removing
tags with fewer than 20 samples, ten classes were retained: Array,
String, Math, Hash Table, Database, Tree, Dynamic Programming, Linked
List, Two Pointers, and Depth First Search. The class distribution
before and after filtering is shown in Figure
\ref{fig:tag-class-distribution}.

\begin{figure}[H]
    \centering
    \includegraphics[width=\linewidth]{./images/evaluation/01_class_distribution.png}
    \caption{LeetCode topic tag distribution before filtering
    (left) and after applying the 60 sample minimum threshold (right).
    Array dominates with 1,168 examples. Tags such as Stack (13),
    Binary Search (10), and Graph (8) were removed due to insufficient
    samples.}
    \label{fig:tag-class-distribution}
\end{figure}

The dataset is substantially imbalanced. Array accounts for 1,168 of
the retained problems while Depth First Search contributes only 26.
This imbalance is the dominant factor shaping classifier performance
and is the root cause of the low F1 macro scores observed across all
models \cite{japkowicz2002imbalance}. A model can achieve high accuracy
simply by performing well on Array and String, which together account
for over 45 percent of the dataset, while performing poorly on the
minority classes. This tension between accuracy and macro F1 is
visible throughout the results.

Figure \ref{fig:tag-word-length} shows the word count distribution per
class. Database problems are notably longer on average, which is
consistent with their SQL heavy descriptions. Tree and Linked List
problems tend to be shorter.

\begin{figure}[H]
    \centering
    \includegraphics[width=\linewidth]
        {./images/evaluation/02_text_length_by_class.png}
    \caption{Word count distribution per topic class for the combined
    title and description text. Database problems are consistently
    longer than other classes. The wide variance within classes
    reflects the diverse nature of problems within each topic.}
    \label{fig:tag-word-length}
\end{figure}

\subsection{Feature Analysis}

Figure \ref{fig:tag-top-words} shows the top 30 terms across the
entire corpus after English stop word removal. Terms such as
\texttt{example}, \texttt{input}, \texttt{output}, and \texttt{nums}
appear across all classes and provide little discriminative signal.
The more informative terms are those that appear prominently in only
one or two classes.

\begin{figure}[H]
    \centering
    \includegraphics[width=\linewidth]{./images/evaluation/03_top_words_overall.png}
    \caption{Top 30 terms across the full LeetCode corpus after
    English stop word removal. High frequency generic terms such as
    \texttt{example}, \texttt{input}, and \texttt{output} appear
    across all topic classes and provide limited discriminative value
    for classification.}
    \label{fig:tag-top-words}
\end{figure}

The per class feature analysis from the final SGD Classifier is shown
in Figure \ref{fig:tag-top-features}. This confirms that the model has
learned genuinely meaningful signals: Database problems are
characterised by terms such as \texttt{table}, \texttt{column}, and
\texttt{column type}; Linked List problems by \texttt{linked},
\texttt{head}, and \texttt{nodes list}; Tree problems by \texttt{tree},
\texttt{root}, and \texttt{binary tree}. The features learned for
Dynamic Programming are notably less clean, including general graph
terms such as \texttt{graph} and \texttt{cities}, which explains the
model's weaker performance on that class.

\begin{figure}[H]
    \centering
    \includegraphics[width=\linewidth]{./images/evaluation/12_top_features_per_class.png}
    \caption{Top 12 classifier weight features per topic class for the
    final SGD Classifier. Classes with distinctive vocabulary
    (Database, Linked List, Tree) show clean, semantically meaningful
    features. Dynamic Programming and Hash Table show more overlap
    with other classes, which corresponds directly to their lower F1
    scores.}
    \label{fig:tag-top-features}
\end{figure}

\subsection{Model Comparison}

Seven candidate models were evaluated on the validation set. Model
selection during training was performed using F1 micro, consistent
with the multi label objective and the use of \texttt{GridSearchCV}.
For evaluation and comparison, F1 macro is reported as the primary
metric, as it provides a more informative view of performance across
all classes in the presence of class imbalance
\cite{zhang2014multilabel}. The models correspond directly to the
training configurations defined in the pipeline: logistic regression
with word features, logistic regression with combined word and
character features, an SGD classifier with word features, and a
classifier chain of logistic regressions with combined features. The
results are shown in Figure \ref{fig:tag-f1-comparison}.

\begin{figure}[H]
    \centering
    \includegraphics[width=\linewidth]
        {./images/evaluation/05_model_comparison_f1.png}
    \caption{Validation set F1 macro comparison across the seven candidate
    topic tag classifier models. The SGD classifier achieved the highest
    F1 macro and was therefore selected as the final deployed model.}
    \label{fig:tag-f1-comparison}
\end{figure}

Figure \ref{fig:tag-all-metrics} presents all three metrics (accuracy,
F1 macro, and F1 weighted) for every model side by side, illustrating
the gap between accuracy and macro F1 that results from the class
imbalance.

\begin{figure}[H]
    \centering
    \includegraphics[width=\linewidth]
        {./images/evaluation/09_grouped_metrics.png}
    \caption{Validation accuracy, F1 macro, and F1 weighted for the seven
    candidate topic tag classifier models. The consistent gap between
    accuracy and F1 macro reflects the class imbalance in the LeetCode
    dataset, where high accuracy can be achieved by performing well on the
    dominant classes alone.}
    \label{fig:tag-all-metrics}
\end{figure}

SGD Classifier achieved the highest validation F1 macro at 0.6464 and
was selected as the deployed model. The effect of hyperparameter
tuning via GridSearchCV is shown in Figure \ref{fig:tag-tuning}. For
most models the gain from tuning was marginal, with the exception of
Multinomial NB which improved substantially from 0.615 to 0.785 in
validation accuracy after tuning, suggesting its default alpha
parameter was poorly suited to this corpus.

\begin{figure}[H]
    \centering
    \includegraphics[width=\linewidth]
        {./images/evaluation/06_default_vs_tuned.png}
    \caption{Comparison of default parameters versus GridSearchCV tuned
    parameters for each model on the validation set. The most notable
    improvement is Multinomial NB, which gained 17 percentage points in
    accuracy after tuning.}
    \label{fig:tag-tuning}
\end{figure}

The five fold cross validation F1 macro score distributions are shown
in Figure \ref{fig:tag-cv-distributions}. SGD Classifier ($0.588 \pm
0.019$) and Linear SVC ($0.581 \pm 0.006$) show the most stable
performance. Multinomial NB shows high variance ($0.575 \pm 0.041$),
indicating sensitivity to the particular fold split and making it a
riskier choice for deployment.

\begin{figure}[H]
    \centering
    \includegraphics[width=\linewidth]
        {./images/evaluation/07_cv_score_distributions.png}
    \caption{Cross validation F1 macro score distributions (five fold)
    for all models with best hyperparameters. SGD Classifier and
    Linear SVC show the tightest distributions, indicating the most
    consistent generalisation across folds.}
    \label{fig:tag-cv-distributions}
\end{figure}

\subsection{Final Model Test Set Results}

The SGD Classifier was evaluated on the held out test set. Final
results are summarised in Table \ref{table:tag-results-final}.

\begin{table}[H]
    \centering
    \begin{tabular}{lcc}
        \hline
        \textbf{Metric} & \textbf{Validation} & \textbf{Test Set} \\
        \hline
        Accuracy    & 0.805 & 0.805 \\
        F1 Macro    & 0.646 & 0.622 \\
        F1 Weighted & 0.802 & \\
        \hline
    \end{tabular}
    \caption{Final SGD Classifier performance on validation and test
    sets. The 2.4 percentage point drop in F1 macro from validation
    to test set is modest and consistent with normal generalisation
    behaviour.}
    \label{table:tag-results-final}
\end{table}

The per class F1 scores across all models on both validation and test
sets are shown in Figures \ref{fig:tag-per-class-val} and
\ref{fig:tag-confusion-sgd}.

\begin{figure}[H]
    \centering
    \includegraphics[width=\linewidth]
        {./images/evaluation/10_per_class_f1_heatmap.png}
    \caption{Per class F1 scores for all models on the validation set,
    with the SGD Classifier test set results shown in the final row.
    Array, Database, and Tree are consistently strong across all
    models. Dynamic Programming, Hash Table, and Two Pointers are
    consistently the weakest classes.}
    \label{fig:tag-per-class-val}
\end{figure}

\begin{figure}[H]
    \centering
    \includegraphics[width=\linewidth]
        {./images/evaluation/11_confusion_matrix_best.png}
    \caption{Row normalised confusion matrix for the SGD Classifier on
    the test set. Array (0.90), Database (1.00), Math (0.86), and Tree
    (1.00) show strong diagonal dominance. Dynamic Programming (0.40)
    and Two Pointers (0.43) show the highest off diagonal confusion,
    primarily misclassified as Array.}
    \label{fig:tag-confusion-sgd}
\end{figure}

The per class results for the final deployed SGD Classifier are
presented in Table \ref{table:tag-per-class}.

\begin{table}[H]
    \centering
    \begin{tabular}{lcc}
        \hline
        \textbf{Topic Class} & \textbf{Test F1} & \textbf{Primary Confusion} \\
        \hline
        Array              & 0.91 & \\
        Database           & 1.00 & \\
        Tree               & 0.83 & \\
        Math               & 0.79 & \\
        Linked List        & 0.71 & Hash Table \\
        Depth First Search & 0.73 & Array \\
        String             & 0.57 & Array \\
        Two Pointers       & 0.29 & String or Math \\
        Dynamic Programming & 0.35 & Array \\
        Hash Table         & 0.20 & Array or String \\
        \hline
    \end{tabular}
    \caption{Per class F1 scores for the final SGD Classifier on the
    test set, with the primary misclassification direction noted for
    weaker classes.}
    \label{table:tag-per-class}
\end{table}

\subsection{Critical Analysis: Class Imbalance and Conceptual Overlap}

Two factors dominate the weaknesses of the topic tag classifier. The
first is class imbalance: Array accounts for 1,168 training examples
while Depth First Search has only 26. This means the classifier has
seen roughly 45 times as many Array examples, and when a minority
class problem contains any vocabulary that also appears in Array
problems (which is almost always true for terms like \texttt{nums},
\texttt{input}, and \texttt{length}), the model is biased towards
predicting Array.

The second factor is conceptual overlap between classes. Dynamic
Programming and Depth First Search both involve graph like reasoning,
and their problem descriptions frequently share terms. Hash Table
problems often involve sliding window or two pointer counting patterns
that resemble both Array and Two Pointers problems. These overlaps are
not an artefact of the model's weakness. They reflect genuine ambiguity
in the label definitions themselves, which is consistent with the
broader observation by Zhang and Zhou \cite{zhang2014multilabel} that
label correlation is a defining feature of multi label problems.

Training time for all models was extremely fast, as shown in
Figure \ref{fig:tag-training-time}, with all models completing
GridSearchCV in under seven seconds. This confirms that the
computational cost of maintaining and retraining the classifier is
negligible.

\begin{figure}[H]
    \centering
    \includegraphics[width=0.8\linewidth]
        {./images/evaluation/08_training_time.png}
    \caption{Total training time including GridSearchCV for each topic
    tag classifier model. All models complete training in under seven
    seconds, confirming that retraining is computationally trivial
    should the dataset be expanded in future.}
    \label{fig:tag-training-time}
\end{figure}

Despite these limitations, a test set accuracy of 0.805 and F1 macro
of 0.622 represents a genuinely useful classification capability for
the platform's purpose. A user generating questions about Dynamic
Programming may occasionally receive a question tagged as Array, but
the majority of generated questions will receive a meaningful and
correct topic label.

\section{Application Level Evaluation}
\label{sec:eval-application}

\subsection{Objective 1: AI Generation Pipeline}

The end to end quiz generation pipeline was tested manually across
both AI backends and both quiz types. Both the OpenAI
(\texttt{gpt-4o-mini}) and Ollama (\texttt{llama3.1:8b}) backends
successfully produced well formed MCQ and coding challenge JSON in
the vast majority of test runs. The OpenAI backend consistently
produced clean, schema conforming JSON that parsed without
intervention. The Ollama backend required additional post processing
in \texttt{parser\_ollama.py}, most notably stripping option prefix
characters such as \texttt{"A: "} from MCQ option strings, which the
OpenAI model did not include.

The most significant challenge encountered during development was
JSON parsing reliability. LLMs do not guarantee valid JSON output,
and both backends occasionally produced responses with minor
formatting errors such as trailing commas, truncated option strings,
or inconsistent field names when the prompt asked for a larger number
of questions. This was partially mitigated by the strict schema
specified in the prompt templates and by wrapping the parse step in
try and except blocks that return meaningful error messages rather
than crashing silently. This parsing fragility is a known limitation
of prompt based LLM integration \cite{liu2023promptsurvey} and is
discussed further in the conclusion. The availability of structured
output enforcement at the API level \cite{openai2024structured}
offers a concrete path to resolving this limitation in a future
iteration.

The 120 second \texttt{httpx} timeout was appropriate in practice.
Requests to \texttt{gpt-4o-mini} for five questions typically
completed in 5 to 15 seconds. Requests to the local Ollama backend
were slower at 20 to 60 seconds depending on hardware, but remained
within the timeout window.

\subsection{Objective 5: Code Execution Engine}

The code execution engine was tested manually across all five
supported languages (Python, JavaScript, TypeScript, Java, C\#) using
both correct and intentionally incorrect solutions. The five second
timeout correctly terminated infinite loop submissions in all tested
cases. Python and JavaScript execution was reliable and fast.
TypeScript required the transpile then execute step which added a
small overhead but functioned correctly. Java and C\# required careful
handling of the class name and entry point expectations of their
respective runtimes, and both worked reliably once these constraints
were enforced in the execution logic.

Test case pass and fail evaluation was verified by submitting known
correct solutions against the AI generated test cases and confirming
that all cases passed, and then introducing deliberate off by one
errors and confirming that the appropriate subset of test cases
failed. The output comparison logic correctly handles trailing
whitespace differences.

\subsection{Objective 6: User Features}

The authentication system, experience and levelling mechanic, quiz
result persistence, and social friends network were all tested through
manual interaction with the running application. JWT protected routes
correctly rejected requests without a valid token. The levelling
formula (\texttt{xp\_for\_level = 100 * level}) produced correct
behaviour across level boundaries in manual testing.

The friends system was tested by creating two user accounts and
exercising the full lifecycle: sending a friend request, accepting it,
verifying mutual friendship, and then removing the friendship. The
\texttt{UNIQUE(requester\_id, addressee\_id)} database constraint
correctly prevented duplicate requests, and the \path{CHECK
(requester\_id != addressee\_id)} constraint correctly rejected self
friendship attempts.

\section{Evaluation Against Objectives}
\label{sec:eval-objectives}

Table \ref{table:objectives-summary} summarises the outcome for each
of the six objectives defined in the Introduction.

\begin{table}[H]
    \centering
    \renewcommand{\arraystretch}{1.4}
    \begin{tabular}{p{0.05\linewidth} p{0.45\linewidth} p{0.12\linewidth}
        p{0.30\linewidth}}
        \hline
        \textbf{\#} & \textbf{Objective} & \textbf{Status} &
        \textbf{Notes} \\
        \hline
        1 & Full stack web application with AI powered MCQ and coding
            challenge generation supporting multiple AI backends
          & Met
          & Both OpenAI and Ollama backends functional. JSON parsing
            occasionally fragile on large requests. \\
        2 & Multi label topic tag classifier predicting algorithmic
            topic tags for generated questions
          & Met
          & Test accuracy 0.805, F1 macro 0.622. Class imbalance limits
            performance on minority classes. \\
        3 & Difficulty classifier predicting easy, medium, or hard in
            the absence of explicit human labels
          & Met
          & XGBoost test accuracy approx. 0.593. Proxy label ceiling
            acknowledged. Medium class hardest to classify. \\
        4 & Both classifiers integrated into the backend inference
            pipeline, enriching every generated question automatically
          & Met
          & Both models loaded as in process singletons. Inference
            adds negligible latency per question. \\
        5 & Code execution engine supporting multiple programming
            languages with test case evaluation
          & Met
          & All five languages functional. Five second timeout
            enforced. No kernel level sandboxing. \\
        6 & User accounts system with authentication, experience and
            levelling, and social friends network
          & Met
          & JWT auth, bcrypt passwords, full friends lifecycle, and
            experience system all functional. \\
        \hline
    \end{tabular}
    \caption{Summary of project objectives and their evaluation
    outcomes.}
    \label{table:objectives-summary}
\end{table}

All six objectives were met. The primary areas of limitation are the
difficulty classifier's reliance on a proxy label, the topic
classifier's performance on minority classes, and the JSON parsing
fragility of the AI generation pipeline. These limitations are not
failures of implementation but reflect genuine challenges in the
problem domain, and each is discussed further in the Conclusion
chapter.